\subsection{IMAGES OF A UNION AND AN INTERSECTION}

PROPOSITION 3. Let \((\mathbf{X}_t)_{t\in I}\) be a family of subsets of a set A, and let \(\Gamma\) be a correspondence between A and B. Then

\[
\Gamma \left\langle \bigcup_ {i \in I} X _ {i} \right\rangle = \bigcup_ {i \in I} \Gamma \left\langle X _ {i} \right\rangle , \quad \Gamma \left\langle \bigcap_ {i \in I} X _ {i} \right\rangle \subset \bigcap_ {i \in I} \Gamma \left\langle X _ {i} \right\rangle .
\]

The relation \((\exists x)\left(x\in \bigcup_{i\in I}X_i \text{ and } y\in \Gamma (x)\right)\) is equivalent to

\[
(\exists x) (\exists t) (t \in I \text {   and   } x \in X _ {t} \text {   and   } y \in \Gamma (x)),
\]

hence to \((\exists \iota)(\iota \in I \text{ and } y \in \Gamma\langle X_{\iota}\rangle)\) , hence is equivalent to \(y \in \bigcup_{\iota \in I} \Gamma\langle X_{\iota}\rangle\) ; this proves the first formula. As to the second formula, we have \(\bigcap_{\iota \in I} X_{\iota} \subset X_{\iota}\) for all \(\iota \in I\) , whence (§3, Proposition 2)

\[
\Gamma \left\langle \bigcap_ {i \in I} X _ {i} \right\rangle \subset \Gamma \left\langle X _ {i} \right\rangle ,
\]

and consequently

\[
\Gamma \left\langle \bigcap_ {i \in I} X _ {i} \right\rangle \subset \bigcap_ {i \in I} \Gamma \left\langle X _ {i} \right\rangle .
\]

If \(\Gamma\) is an arbitrary correspondence (and in particular an arbitrary function), the formula

\[
\Gamma \left\langle \bigcap_ {i \in I} X _ {i} \right\rangle = \bigcap_ {i \in I} \Gamma \left\langle X _ {i} \right\rangle
\]

is usually false.

* For example, in the plane \(\mathbf{R}^2\) the first projections of the lines \(y = x\) and \(y = x + 1\) are equal to \(\mathbf{R}\) , but the intersection of these lines is empty, and therefore so is the first projection of this intersection (*). *

However, we have the following important result:

PROPOSITION 4. Let \(f\) be a mapping of A into B and let \((\mathbf{Y}_t)_{t \in I}\) be a family of subsets of B. Then \(\overline{f}^1 \left\langle \bigcap_{i \in I} \mathbf{Y}_t \right\rangle = \bigcap_{i \in I} \overline{f}^1 \langle \mathbf{Y}_t \rangle\) .

To prove this, let \(x\) be an element of \(\bigcap_{i \in I} f^{-1} \langle Y_i \rangle\) . We have \(f(x) \in Y_i\) for all \(i \in I\) , whence \(f(x) \in \bigcap_{i \in I} Y_i\) , and consequently

\[
x \in \overline {{f}} ^ {1} \left\langle \bigcap_ {\iota \in I} Y _ {\iota} \right\rangle .
\]

Therefore \(\bigcap_{i\in I}\overline{f}^{1}(Y_{i})\subset\overline{f}^{1}\Big\langle\bigcap_{i\in I}Y_{i}\Big\rangle\) ; this relation, together with Proposition 3, gives the result.

COROLLARY. If \(f\) is an injection of \(A\) into \(B\) and if \((X_{t})_{t \in I}\) is a family of subsets of \(A\) whose index set \(I\) is not empty, then \(f\left\langle \bigcap_{i \in I} X_{t} \right\rangle = \bigcap_{i \in I} f\langle X_{t} \rangle\) .

For we may write \(f = i \circ g\) , where \(i\) is the canonical injection of \(f\langle \mathbf{A} \rangle\) into \(\mathbf{B}\) and \(g\) is a bijection of \(\mathbf{A}\) onto \(f\langle \mathbf{A} \rangle\) . If \(h\) denotes the inverse mapping of \(g\) , we have \(f\langle \mathbf{X} \rangle = \overline{h}^1\langle \mathbf{X} \rangle\) for every subset \(\mathbf{X}\) of \(\mathbf{A}\) , and we are therefore brought back to Proposition 4.

